\chapter{Introduction: three games and one exponent}\label{ch:intro}

Goldbach's conjecture asserts that every even integer $\N>2$ is a sum of two
primes; it has been verified for all $\N\le4\cdot10^{18}$~\cite{OeS}.
Heuristically one expects not just one representation but many. The
Hardy--Littlewood circle method---which estimates such counts by treating the
primes as a random set of the appropriate density and then correcting for
divisibility constraints---predicts~\cite{HL}
\begin{equation}\label{eq:HL}
r(\N)\;\sim\;2\,\Pi_2\,\frac{\N}{(\log \N)^2}\prod_{\substack{p\mid \N\\ p>2}}\frac{p-1}{p-2},
\qquad
\Pi_2=\prod_{p>2}\Bigl(1-\tfrac{1}{(p-1)^2}\Bigr),
\end{equation}
where $r(\N)=\#\{p\le\N:p\text{ and }\N-p\text{ both prime}\}$. Reading this
from left to right: $\N/(\log\N)^2$ is the naive guess, since a number near
$\N$ is prime with probability about $1/\log\N$ and there are $\N$ choices for
$p$; the twin-prime constant $\Pi_2\approx0.6601$ corrects for the fact that
primes are not truly random modulo small numbers; and the last factor inflates
the count when $\N$ is highly composite. The conclusion is that a typical large
even integer should have \emph{many} Goldbach partitions. So the interesting
question is not whether a representation exists, but how lopsided the most
lopsided one must be.

\begin{definition}
For an even integer $\N\ge4$ admitting a Goldbach representation, the
\emph{least Goldbach summand} is
\[
\g(\N)\;=\;\min\{\,p\ \text{prime}:\N-p\ \text{is prime}\,\}.
\]
\end{definition}

This is the quantity $p(n)$ of the \emph{minimal Goldbach partition} studied by
Granville, van de Lune and te Riele~\cite{GLvdtR}, who conjectured a ceiling
\begin{equation}\label{eq:GLR}
\g(\N)\;=\;O\bigl((\log \N)^2\log\log \N\bigr).
\end{equation}
Empirically $\g$ is tiny and grows very slowly; its record values are
catalogued in the OEIS~\cite{OEIS} (A020481, with records A025018/A025019).
Exhaustive search gives a slowly climbing ladder: $\g=1093$ at
$\N=60\,119\,912$; $\g(\N)\le5569$ for all even
$\N\le4\cdot10^{14}$~\cite{Richstein}; and the current exhaustive record
$\g=9781$, first attained at $\N=3\,325\,581\,707\,333\,960\,528$, with
$\g(\N)\le9781$ for \emph{all} even $\N\le4\cdot10^{18}$~\cite{OeS}. Strikingly,
these records occur at integers of at most $19$ digits.

We ask the opposite, extremal and constructive question:
\begin{quote}\itshape
How large can $\g(\N)$ be forced to be, and at what cost in the size of $\N$?
\end{quote}
An even integer $\N$ for which $\g(\N)$ is large is a \emph{Goldbach desert}:
a number all of whose small prime offsets are blocked. Constructing deserts is
the subject of this thesis, and the reason the subject is interesting is that
every construction turns out to be governed by a single number---an exponent
that measures how much luck the construction still needs after all its
cleverness is spent.

\section{The three games}

There are three natural record questions, and they reward different
constructions. The first is the unbounded \emph{height game}: exhibit as large
a value of $\g(\N)$ as possible, with no limit on the size of $\N$. An ordinary
prime gap is one way to play, but it is a strictly stronger construction: a gap
must block \emph{every} integer offset, while a desert needs to block only the
\emph{prime} offsets. Prime-gap constructions are therefore a strict subclass of
the constructions available here---a point we return to in
\S\ref{sec:gaps} with a concrete two-orders-of-magnitude example.

For comparisons in which size matters we use two finite games.

\begin{definition}[Threshold and budget games]\label{def:games}
For $Q\ge2$ and $X\ge4$, define
\[
T(Q)=\min\{\,\N \text{ even, Goldbach-representable}:\ \g(\N)>Q\,\},
\qquad
R(X)=\max\{\,\g(\N):\ \N<X\,\}.
\]
The \emph{threshold game} fixes $Q$ and tries to lower $T(Q)$; the
\emph{budget game} fixes $X$ and tries to raise $R(X)$. ``Fewer than $200$
decimal digits'' means the strict budget $\N<10^{199}$.
\end{definition}

\begin{remark}[A naming caution]
The budget game has been written both as $R(10^{199})$ (``fewer than $200$
digits'') and as $R(10^{200})$ (``at most $200$ digits'') in the literature
this thesis builds on, and the two are genuinely different: our $200$-digit
incumbent with $\g=112\,249$ satisfies $\N<10^{200}$ but exceeds $10^{199}$ by
$3.4\%$, which is precisely why a $199$-digit integer with the smaller value
$\g=119\,419$ is the record for the stricter budget. Throughout we state the
bound explicitly with each record.
\end{remark}

\section{Certified incumbents}

Table~\ref{tab:records} lists the programme's certified incumbents at the time
of writing. Every entry has been verified twice: once by the pipeline that
produced it and once by an independent verifier written against the record file
rather than against the search code, including a re-derivation of every
per-offset witness and an independent check of the elliptic-curve primality
certificate.

\begin{table}[h]
\centering\small
\caption{Certified incumbents. ``Witnesses'' counts prime offsets excluded by
an explicit congruence divisor plus those excluded by a strong base-$2$
compositeness witness; in every case the positive side is an exhibited prime
complement with an ECPP certificate.}\label{tab:records}
\begin{tabular}{lrrl}
\toprule
Game & Digits of $\N$ & $\g(\N)$ & Construction\\
\midrule
Height (unbounded) & $2\,480$ & $1\,157\,341$ & cover at $Q=1.16\cdot10^6$, $t=14\,544$\\
Height, gap-dominating & $2\,692$ & $1\,134\,871$ & cover at $Q=1\,113\,138$, $k=14$\\
Threshold $T(100\,000)$ & $150$ & $104\,527$ & cover at $Q=100\,003$, $k=6.21\cdot10^{10}$\\
Budget $R(10^{199})$ & $199$ & $119\,419$ & cover at $Q\approx1.19\cdot10^5$, $t=1.46\cdot10^8$\\
Budget $R(10^{200})$ & $200$ & $112\,249$ & cover at $Q=107\,720$, $k=56\,770$\\
\bottomrule
\end{tabular}
\end{table}

Two of these deserve comment. The $2\,692$-digit entry is included because it
settles the comparison with ordinary prime gaps in the strongest possible way:
its $\g$ exceeds the largest \emph{certified} prime gap known
($1\,113\,106$), and it does so with an integer roughly $6.9$ times shorter
than the one realising that gap. Deserts really are cheaper than gaps, by a
measurable factor.

The $150$-digit threshold record is the endpoint of a \emph{ratchet}: the same
covering machinery, re-tuned and re-run, produced a descending ladder
$195\to179\to150$ digits at $Q\approx10^5$, banking a certified record at each
rung. Chapter~\ref{ch:engineering} argues that this ratchet discipline, rather
than any single clever cover, is the strategy that actually produces records
under a fixed compute budget.

\section{The exponent that governs everything}

The technical spine of the thesis is a single quantity. A partial congruence
cover fixes $\N$ inside an arithmetic progression $\N=\N_0+kM$ and disposes of
most prime offsets for \emph{every} $k$; a residual set $U$ of offsets survives
and must be knocked out by luck, one candidate $k$ at a time. Under the
standard heuristic, the probability that a given candidate succeeds is
$e^{-E}$ with
\begin{equation}\label{eq:E}
E\;=\;\frac{|U|\cdot\mathrm{boost}}{\log\N},
\qquad
\mathrm{boost}\;=\;2\prod_{r\in\mathcal R}\frac{r}{r-1},
\end{equation}
so that $e^{E}$ candidates must be examined per expected hit. We call $E$ the
\emph{failure exponent}. It is the currency of the entire subject: every design
decision---which moduli to use, which residues, how many offsets to leave
residual, what digit budget to spend---is a trade in units of $E$, and every
piece of engineering (faster tests, better sieving, selective testing, more
hardware) buys a fixed number of nats of $E$ and no more.

Equation~\eqref{eq:E} is a heuristic, so we test it rather than assume it.
Across three independent covers and $1.2\times10^{9}$ scanned candidates, the
measured exponent $\Ehat$ agrees with the predicted $E$ to within $0.1$ nats
(Chapter~\ref{ch:engineering}). That agreement is what licenses the rest of the
thesis: once the model is trusted, questions about feasibility become
arithmetic.

\section{What is new here}

Chapters~\ref{ch:foundations} and~\ref{ch:engineering} are largely expository
and engineering: they set out the classical constructions, the covering
formulation, the calibrated cost model, and the search engine, and they report
the certified records and---equally important---the campaigns that failed.
The contributions begin in Chapter~\ref{ch:frontier}.

\begin{enumerate}
\item \textbf{The difficulty is an integrality gap} (\S\ref{sec:lp}). The
natural linear-programming relaxation of the covering problem is
\emph{vacuous}: fractionally, the available moduli cover $100.00\%$ of the
prime offsets at every target we tested. Consequently the whole failure
exponent is the price of \emph{rounding}: naive independent rounding leaves
$15.6\%$ of offsets uncovered, adaptive greedy leaves $7.56\%$, and the
fractional optimum leaves none. Every absolute percentage point recovered is
worth about $2.7$ nats, i.e.\ a factor of $15$ in compute.

\item \textbf{The rounding gap resists optimisation} (\S\ref{sec:door1}). We
attack it with exact large-neighbourhood search (certified optimal by LP
duality on $21$ of $22$ windows), with belief propagation and decimation, and
with a simultaneous exact solve of half the cover at once. The greedy solution
$|U|=867$ is never beaten. A random-restart probe shows why: the landscape is a
field of mutually distant, locally-exact-stable basins, and the adaptive greedy
construction lands about $100$ offsets deeper than a typical one.

\item \textbf{Sub-$100$-digit deserts reduce to a hard problem}
(Chapter~\ref{ch:barriers}). Oversized covers would make sub-$100$-digit
deserts easy if their Chinese-remainder representatives could be made small;
that requirement is exactly a modular subset-sum problem with per-position
choices at density $\approx1.026$, the regime where neither lattice reduction
nor birthday methods work. We also verify against the source that the relevant
CRT list-decoding radius is a factor $5.2$ short of what the construction would
need, and confirm experimentally that lattice methods do not penetrate the
Johnson-type radius even when solutions are abundant.

\item \textbf{A closure theorem} (Chapter~\ref{ch:closure}). Within a concrete
language of certificate schemas, no construction beats covering. Every family
closes for its own quantifiable reason---a $\sqrt{Q}$ value-set ceiling, a
purchase plan running $2.5\times$ over the design-entropy budget, a cluster
bound, lattice isolation, or an empty intersection at record scale---and the
theorem is conditional on two named hypotheses, one of which we pose as an open
problem.

\item \textbf{A methodological finding} (\S\ref{sec:referee}). We stress-tested
the closure theorem with an adversarial evolutionary search against a
machine-checkable scorer. Its first success was against the \emph{referee}, not
the theorem: it found a $50\times$ accounting error in our evaluator and
exploited it within ten candidates. The audit protocol caught it before any
claim was made, and repairing the referee strengthened the mathematics.
\end{enumerate}

\section{Notation and conventions}

Throughout, $\log_k$ is the $k$-fold iterated logarithm (so
$\log_2\N=\log\log\N$), and each extra $\log$ grows \emph{extremely} slowly;
$\theta(x)=\sum_{p\le x}\log p$ is Chebyshev's function;
$\gamma\approx0.5772$ is Euler's constant; $f\gg g$ means $f\ge cg$ for some
$c>0$ and all large arguments; $\pi(x)$ is the prime-counting function; and
``BPSW'' denotes the Baillie--PSW probable-prime test~\cite{PSW,BW}. All
entropies, costs and exponents are measured in \emph{nats} (natural logarithm
units), because the natural unit of a congruence modulo $r$ is $\log r$ and the
natural unit of a search is $\log(\text{candidates})$.

We are careful to separate three claims about any proposed record, and the
reader should hold us to the distinction throughout:
\begin{itemize}
\item[(a)] no prime below the stated value is a Goldbach summand of $\N$;
\item[(b)] $\N$ admits a Goldbach representation at all;
\item[(c)] the least summand \emph{equals} the stated value.
\end{itemize}
A congruence or a failed probable-prime test settles (a) unconditionally---a
prime can never fail such a test, so negative results are never in doubt. Claim
(b) needs an exhibited prime complement, and (c) needs both together with a
primality certificate for that complement. Where a statement is heuristic,
conditional, or merely measured, we say so: results proved at the level of
rigour of this thesis are marked as such, statements resting on standard but
unproved analytic heuristics are flagged, and purely empirical claims are
reported with the size of the experiment that supports them.
